\documentclass[10pt,a4paper]{amsart}
\usepackage{amsmath,amssymb,amsthm,graphicx}
\usepackage{verbatim}
\newenvironment{reference}{\comment}{\endcomment}
%\newenvironment{reference}{}{}
\newenvironment{slogan}{\comment}{\endcomment}
\newenvironment{history}{\comment}{\endcomment}

% For commutative diagrams we use Xy-pic
\usepackage[all]{xy}

% We use 2cell for 2-commutative diagrams.
\xyoption{2cell}
% Theorem environments.
%
\theoremstyle{plain}
\newtheorem{theorem}[subsection]{Theorem}
\newtheorem{proposition}[subsection]{Proposition}
\newtheorem{lemma}[subsection]{Lemma}

\theoremstyle{definition}
\newtheorem{definition}[subsection]{Definition}
\newtheorem{example}[subsection]{Example}
\newtheorem{exercise}[subsection]{Exercise}
\newtheorem{situation}[subsection]{Situation}

\theoremstyle{remark}
\newtheorem{remark}[subsection]{Remark}
\newtheorem{remarks}[subsection]{Remarks}

\numberwithin{equation}{subsection}

% Macros
%
\def\lim{\mathop{\mathrm{lim}}\nolimits}
\def\colim{\mathop{\mathrm{colim}}\nolimits}
\def\Spec{\mathop{\mathrm{Spec}}}
\def\Hom{\mathop{\mathrm{Hom}}\nolimits}
\def\Ext{\mathop{\mathrm{Ext}}\nolimits}
\def\SheafHom{\mathop{\mathcal{H}\!\mathit{om}}\nolimits}
\def\SheafExt{\mathop{\mathcal{E}\!\mathit{xt}}\nolimits}
\def\Sch{\mathit{Sch}}
\def\Mor{\mathop{\mathrm{Mor}}\nolimits}
\def\Ob{\mathop{\mathrm{Ob}}\nolimits}
\def\Sh{\mathop{\mathit{Sh}}\nolimits}
\def\NL{\mathop{N\!L}\nolimits}
\def\CH{\mathop{\mathrm{CH}}\nolimits}
\def\proetale{{pro\text{-}\acute{e}tale}}
\def\etale{{\acute{e}tale}}
\def\QCoh{\mathit{QCoh}}
\def\Ker{\mathop{\mathrm{Ker}}}
\def\Im{\mathop{\mathrm{Im}}}
\def\Coker{\mathop{\mathrm{Coker}}}
\def\Coim{\mathop{\mathrm{Coim}}}

% Boxtimes
%
\DeclareMathSymbol{\boxtimes}{\mathbin}{AMSa}{"02}

%
% Macros for moduli stacks/spaces
%
\def\QCohstack{\mathcal{QC}\!\mathit{oh}}
\def\Cohstack{\mathcal{C}\!\mathit{oh}}
\def\Spacesstack{\mathcal{S}\!\mathit{paces}}
\def\Quotfunctor{\mathrm{Quot}}
\def\Hilbfunctor{\mathrm{Hilb}}
\def\Curvesstack{\mathcal{C}\!\mathit{urves}}
\def\Polarizedstack{\mathcal{P}\!\mathit{olarized}}
\def\Complexesstack{\mathcal{C}\!\mathit{omplexes}}
% \Pic is the operator that assigns to X its picard group, usage \Pic(X)
% \Picardstack_{X/B} denotes the Picard stack of X over B
% \Picardfunctor_{X/B} denotes the Picard functor of X over B
\def\Pic{\mathop{\mathrm{Pic}}\nolimits}
\def\Picardstack{\mathcal{P}\!\mathit{ic}}
\def\Picardfunctor{\mathrm{Pic}}
\def\Deformationcategory{\mathcal{D}\!\mathit{ef}}
\usepackage{cleveref}
\crefname{originalequation}{original equation}{original equations}
\Crefname{originalequation}{Original equation}{Original equations}
\makeatletter
\newlabel{equation-FM-QCoh}{{0FZC}{source}{}{sourcepointer-1}{}}
\newlabel{equation-FM-QCoh@cref}{{[originalequation][0][]0FZC}{[1][0][]source}{}{}{}}
\makeatother
\crefname{originallemma}{original lemma}{original lemmas}
\Crefname{originallemma}{Original lemma}{Original lemmas}
\makeatletter
\newlabel{lemma-functor-quasi-coherent-from-affine-diagonal-pre}{{0GPA}{source}{}{sourcepointer-2}{}}
\newlabel{lemma-functor-quasi-coherent-from-affine-diagonal-pre@cref}{{[originallemma][0][]0GPA}{[1][0][]source}{}{}{}}
\makeatother
\crefname{originallemma}{original lemma}{original lemmas}
\Crefname{originallemma}{Original lemma}{Original lemmas}
\makeatletter
\newlabel{lemma-coh-noetherian-from-affine-flat}{{0FZG}{source}{}{sourcepointer-3}{}}
\newlabel{lemma-coh-noetherian-from-affine-flat@cref}{{[originallemma][0][]0FZG}{[1][0][]source}{}{}{}}
\makeatother
\title{Exact functors over quasi-compact schemes with affine diagonal are exactly the FM functors with flat acyclic kernels}
\author{The Stacks Project Authors}
\date{stacks-project:functors.tex:local-snapshot-20261012; GNU FDL 1.2 or later}
\begin{document}
\maketitle
\noindent\textit{Editorial note.} All delivered mathematical text is the author's own, copied verbatim from the pinned Stacks Project chapter file functors.tex; no statement or proof step was rewritten, reordered or removed.\par\medskip
\noindent\textit{Editorial note.} Span recomputation by label pairing (LANE\_BRIEF 10.2): in this chapter the author's proof environment follows the statement environment as a standalone block; the spec covers the complete environment pair (statement span, proof span) as recorded in the pool candidate.\par\medskip
\noindent\textit{Editorial note.} Occupation scan: combined claim\_key + statement\_anchor + generic-line scan over all registration files (active and withdrawn) found 0 hits for functors.tex (fu-occupied.json, 2026-10-14); consistent with the duty manager's six-check dispatch.\par\medskip
\ifdefined\hypertarget\hypertarget{sourcepointer-1}{}\fi
\noindent\textit{Source reference.} Original equation 0FZC: sources/stacks-chow/tags.txt, source lines 19947--19947. Cross-chapter standard prerequisite cited by the author inside the counted span; it is used as a pointer and its own proof is not part of this item.\par\medskip
\ifdefined\hypertarget\hypertarget{sourcepointer-2}{}\fi
\noindent\textit{Source reference.} Original lemma 0GPA: sources/stacks-chow/tags.txt, source lines 20820--20820. Cross-chapter standard prerequisite cited by the author inside the counted span; it is used as a pointer and its own proof is not part of this item.\par\medskip
\ifdefined\hypertarget\hypertarget{sourcepointer-3}{}\fi
\noindent\textit{Source reference.} Original lemma 0FZG: sources/stacks-chow/tags.txt, source lines 19951--19951. Cross-chapter standard prerequisite cited by the author inside the counted span; it is used as a pointer and its own proof is not part of this item.\par\medskip
\begin{lemma}
\label{lemma-functor-quasi-coherent-from-affine-diagonal}
Let $R$ be a ring. Let $X$ and $Y$ be schemes over $R$. Assume $X$ is
quasi-compact with affine diagonal.
There is an equivalence of categories between
\begin{enumerate}
\item the category of $R$-linear exact functors
$F : \QCoh(\mathcal{O}_X) \to \QCoh(\mathcal{O}_Y)$
which commute with arbitrary direct sums, and
\item the full subcategory of $\QCoh(\mathcal{O}_{X \times_R Y})$ consisting
of $\mathcal{K}$ such that
\begin{enumerate}
\item $\mathcal{K}$ is flat over $X$,
\item for $\mathcal{F} \in \QCoh(\mathcal{O}_X)$ we have
$R^q\text{pr}_{2, *}(\text{pr}_1^*\mathcal{F}
\otimes_{\mathcal{O}_{X \times_R Y}} \mathcal{K}) = 0$ for $q > 0$.
\end{enumerate}
\end{enumerate}
given by sending $\mathcal{K}$ to the functor $F$ in (\ref{equation-FM-QCoh}).
\end{lemma}

\begin{proof}
Let $\mathcal{K}$ be as in (2). The functor $F$ in
(\ref{equation-FM-QCoh}) commutes with direct sums.
Since by (1) (a) the modules $\mathcal{K}$ is $X$-flat,
we see that given a short exact
sequence $0 \to \mathcal{F}_1 \to \mathcal{F}_2 \to \mathcal{F}_3 \to 0$
we obtain a short exact sequence
$$
0 \to
\text{pr}_1^*\mathcal{F}_1 \otimes_{\mathcal{O}_{X \times_R Y}} \mathcal{K} \to
\text{pr}_1^*\mathcal{F}_2 \otimes_{\mathcal{O}_{X \times_R Y}} \mathcal{K} \to
\text{pr}_1^*\mathcal{F}_3 \otimes_{\mathcal{O}_{X \times_R Y}} \mathcal{K} \to
0
$$
Since by (2)(b) the higher direct image $R^1\text{pr}_{2, *}$
on the first term is zero, we conclude that
$0 \to F(\mathcal{F}_1) \to F(\mathcal{F}_2) \to F(\mathcal{F}_3) \to 0$
is exact and we see that $F$ is as in (1).

\medskip\noindent
Let $F$ be as in (1). Let $\mathcal{K}$ and $t : F \to F_\mathcal{K}$ be as
in Lemma \ref{lemma-functor-quasi-coherent-from-affine-diagonal-pre}.
By Lemma \ref{lemma-coh-noetherian-from-affine-flat} we see
that $\mathcal{K}$ is flat over $X$. To finish the proof we have
to show that $t$ is an isomorphism and the statement on higher
direct images. Both of these follow from the fact that the
relative {\v C}ech complex
$$
\bigoplus 
\text{pr}_{i_0, 2, *}
(\text{pr}_1^*\mathcal{F} \otimes \mathcal{K})_{i_0}
\to
\bigoplus 
\text{pr}_{i_0i_1, 2, *}
(\text{pr}_1^*\mathcal{F} \otimes \mathcal{K})_{i_0i_1}
\to
\bigoplus 
\text{pr}_{i_0i_1i_2, 2, *}
(\text{pr}_1^*\mathcal{F} \otimes \mathcal{K})_{i_0i_1i_2}
\to \ldots
$$
computes $R\text{pr}_{2, *}(\text{pr}_1^*\mathcal{F} \otimes \mathcal{K})$.
Please see proof of
Lemma \ref{lemma-functor-quasi-coherent-from-affine-diagonal-pre}
for notation and for the reason why this is so. In the proof of
Lemma \ref{lemma-functor-quasi-coherent-from-affine-diagonal-pre}
we also found that this complex is equal to $F$ applied to the complex
$$
\bigoplus j_{i_0*}\mathcal{F}_{i_0} \to
\bigoplus j_{i_0i_1*}\mathcal{F}_{i_0i_1} \to
\bigoplus j_{i_0i_1i_2*}\mathcal{F}_{i_0i_1i_2} \to \ldots
$$
This complex is exact except in degree zero with cohomology
sheaf equal to $\mathcal{F}$. Hence since $F$ is an exact functor
we conclude $F = F_\mathcal{K}$ and that (2)(b) holds.

\medskip\noindent
We omit the proof that the construction that sends $F$ to
$\mathcal{K}$ is functorial and a quasi-inverse to the
functor sending $\mathcal{K}$ to the functor $F_\mathcal{K}$
determined by (\ref{equation-FM-QCoh}).
\end{proof}

\end{document}
